\section{Higgs 真空与电弱破缺}\label{sec:sm-higgs-ewsb-dp11}

手征规范结构在未破缺真空中给出一个\emph{无质量的规范理论}：规范对称性允许 Yukawa 三场算符，却仍禁止直接写 $m_W^2W_\mu W^\mu$、$m_Z^2Z_\mu Z^\mu$ 和费米子裸 Dirac 质量。这里的\emph{真空}不是“所有场都等于零”的同义词，而是量子场论中最低能量的参考态；若某个场算符在这个最低能态中的期望值非零，就称它具有非零真空期望值。这里在选定规范后围绕 Higgs 势的经典最低能构型展开；未经规范固定的完整平均不会为规范非单态场赋予可观测的方向期望值。Higgs 相的物理内容由规范不变的谱、质量与散射振幅表述。

拉格朗日量仍具有完整 $SU(2)_L\times U(1)_Y$ 规范对称性。选定规范及一个非零经典真空代表元以后，保持该代表元的生成元构成稳定子；在这个代表元附近的微扰展开通常称为电弱自发对称性破缺。局域规范冗余仍保留在量子理论中，物理谱中的质量产生并不取消规范一致性条件。普通力学中可把它类比成一圈等高的连续极小值：方程对绕中心旋转对称，但任选一个极小点作为展开中心后，该点本身只被较小的稳定子保持。对标准模型，先求 Higgs 势的真空流形，再把 $|D_\mu\Phi|^2$ 在真空上展开。势能沿真空流形切向没有恢复力的方向对应 Goldstone 坐标；在规范理论中，这三个坐标与三个有质量规范场的纵向自由度重组，而径向涨落保留下来成为物理 Higgs 场。对角化规范场质量矩阵以后，唯一的零本征方向才被识别为电磁 $U(1)$，此时电荷算符 $Q$ 和弱混合角才有明确物理意义。

\subsection{Higgs 真空与径向模}

Higgs 二重态包含四个实场自由度。势能先决定真空模长；选定一个真空方向以后，沿真空轨道切向的三个方向对应被破缺生成元，径向方向则是具有恢复力的物理标量涨落。随后把同一个真空代入协变动能，三个切向自由度与三个有质量规范场的纵向极化结合，而径向模保留下来成为物理 Higgs 场。

\begin{definition}[Higgs 势与真空径向尺度]
对 \eqnrefs{eq:sm-higgs-rep-dp11} 的复二重态 $\Phi$，洛伦兹不变性、$SU(2)_L\times U(1)_Y$ 规范不变性与可重整化性把标量拉氏量限制为
\begin{equation}
 \boxed{
 \Lag_{\rm Higgs}
 =(D_\mu\Phi)^\dagger(D^\mu\Phi)
 -V(\Phi),
 \qquad
 V(\Phi)=-\mu_H^2\Phi^\dagger\Phi
 +\frac{\lambda_H}{\hbar c}(\Phi^\dagger\Phi)^2.}
 \label{eq:sm-higgs-lagrangian-dp11}
\end{equation}
取 $\mu_H^2>0$、$\lambda_H>0$。由于势能只依赖规范不变量 $r^2\equiv\Phi^\dagger\Phi$，真空首先由这个径向变量的极小值决定。非零极小值满足 $r^2=\mu_H^2\hbar c/(2\lambda_H)$，因此可把真空代表元与能标定义为
\begin{equation}
 \boxed{
 \langle\Phi\rangle
 =\frac1{\sqrt2}\begin{pmatrix}0\\v\end{pmatrix},
 \qquad
 v^2=\frac{\mu_H^2\hbar c}{\lambda_H},
 \qquad
 v_E\equiv\sqrt{\hbar c}\,v.}
 \label{eq:sm-higgs-vev-dp11}
\end{equation}
$v_E$ 具有能量单位，是电弱质量与 Higgs 耦合共同使用的真空能标。
\end{definition}

\begin{definition}[真空轨道坐标与幺正规范]
一般场构型可沿规范轨道与径向方向分解为
\begin{equation*}
 \Phi(x)
 =\frac1{\sqrt2}
 U_\xi(x)
 \begin{pmatrix}0\\v+h(x)\end{pmatrix},
 \qquad
 U_\xi(x)=\exp\!\left[
 \frac{\ii\xi^I(x)\sigma^I}{v}
 \right].
\end{equation*}
$\xi^I$ 描述三个规范轨道切向自由度。所谓\emph{幺正规范}（unitary gauge）是利用局域 $SU(2)_L\times U(1)_Y$ 规范自由度，把这三个沿规范轨道的坐标选为零，即取 $U_\xi=I$。在这种规范下，被破缺方向的 Goldstone 场不再显式出现在标量二重态中，它们对应的自由度已经进入 $W^\pm$ 与 $Z$ 的纵向极化；标量扇区只剩径向场 $h$：
\begin{equation}
 \Phi_U(x)=\frac1{\sqrt2}\begin{pmatrix}0\\v+h(x)\end{pmatrix}.
 \label{eq:sm-higgs-unitary-gauge-dp11}
\end{equation}
\end{definition}

\begin{proposition}[径向 Higgs 质量与自相互作用]
真空附近的质量与自相互作用由势能沿径向方向的展开系数决定：
\begin{equation}
 \boxed{
 V(h)=V_0+\mu_H^2h^2
 +\frac{\lambda_Hv}{\hbar c}h^3
 +\frac{\lambda_H}{4\hbar c}h^4.}
 \label{eq:sm-higgs-potential-expanded-dp11}
\end{equation}
二次项的曲率定义物理 Higgs 质量，得到
\begin{equation}
 \boxed{
 m_hc^2=\sqrt{2\lambda_H}\,v_E.}
 \label{eq:sm-higgs-mass-dp11}
\end{equation}
因此 Higgs 质量固定四次耦合常数：
\begin{equation*}
 \lambda_H=\frac{(m_hc^2)^2}{2v_E^2}.
\end{equation*}
\end{proposition}

\begin{figure}[htbp]
 \centering
 \includegraphics[width=.72\textwidth]{sm_higgs_potential_dp11.png}
 \caption{Higgs 势能的径向截面。$\mu_H^2>0$ 时原点是不稳定驻点，连续极小值满足 $\Phi^\dagger\Phi=v^2/2$。规范固定选定一个真空代表元后，径向曲率给出物理 Higgs 质量；规范轨道上的三个切向自由度与有质量弱规范玻色子的纵向偏振重组。}
 \label{fig:sm-higgs-potential-dp11}
\end{figure}

\subsection{电弱质量矩阵}

规范玻色子的质量可以直接读作“真空沿哪个生成元方向发生变化”。先在一般标量表示中
建立这一事实，再代入 Higgs 二重态，便能看清零质量光子为什么必然存在。
\begin{proposition}[规范质量二次型与真空稳定子]
设复标量真空代表元为 $v_\Phi$，厄米内部生成元为 $T^a$，几何规范场为
$\mathcal A_\mu^a$、无量纲耦合为 $g_a$。由 $D=\partial+\ii g_aT^a\mathcal A^a$，
能量平方质量矩阵是
\begin{equation}
(\mathsf M_E^2)_{ab}
=\hbar c\,g_ag_b\,v_\Phi^\dagger\{T^a,T^b\}v_\Phi.
\label{eq:sm-general-gauge-mass-gram-dp126}
\end{equation}
它实对称且半正定。对实向量 $z$，$\mathsf M_E^2z=0$ 当且仅当
$\sum_a z_ag_aT^av_\Phi=0$。
\end{proposition}
\begin{proof}
均匀真空的协变动能给
\begin{align*}
(D_\mu v_\Phi)^\dagger D^\mu v_\Phi
&=\sum_{a,b}g_ag_b\mathcal A_\mu^a\mathcal A^{b\mu}
v_\Phi^\dagger T^aT^bv_\Phi\\
&=\frac1{2\hbar c}\mathcal A_\mu^a(\mathsf M_E^2)_{ab}\mathcal A^{b\mu},\\
z^{\mathsf T}\mathsf M_E^2z
&=2\hbar c\left\|\sum_a z_ag_aT^av_\Phi\right\|^2\ge0.
\end{align*}
第一行中的两个规范场可交换，所以只有生成元的对称乘积留下。
第二式既证明非负性，也识别零空间；正交对角化后非零本征值就是 $(m_Vc^2)^2$。
由于 $[v_\Phi^\dagger v_\Phi]=\mathrm{J/m}$，乘 $\hbar c$ 后矩阵确有能量平方量纲。
\end{proof}

\begin{proof}[证明（电弱二重态的质量二次型）]
Higgs 场属于颜色表示 $\mathbf1$，所以 $SU(3)_c$ 生成元在该表示上作用为零；协变导数中的颜色联络项因此消失，留下电弱规范场：
\begin{equation*}
 D_\mu\Phi
 =\left[
 \partial_\mu
 +\ii g\frac{\sigma^I}{2}\mathcal W_\mu^I
 +\ii g'\frac12\mathcal B_\mu
 \right]\Phi.
\end{equation*}
把真空\eqnrefs{eq:sm-higgs-vev-dp11} 代入 $|D_\mu\Phi|^2$ 后，规范场的二次型在对角化之前统一写成
\begin{align}
 (D_\mu\langle\Phi\rangle)^\dagger(D^\mu\langle\Phi\rangle)
 =&\frac{v^2g^2}{8}\left[(\mathcal W^1)^2+(\mathcal W^2)^2\right]\\
 &+\frac{v^2}{8}(g\mathcal W^3-g'\mathcal B)^2.
 \label{eq:sm-gauge-mass-before-diagonal-dp11}
\end{align}
这个二次型自然分成两个线性代数问题：$(\mathcal W^1,\mathcal W^2)$ 构成带电扇区，而 $(\mathcal W^3,\mathcal B)$ 构成需要对角化的 $2\times2$ 中性质量矩阵。
\end{proof}

\begin{proposition}[带电规范场的质量]
带电扇区由下列线性组合对角化：
\begin{equation}
 \boxed{
 \mathcal W_\mu^\pm
 =\frac{\mathcal W_\mu^1\mp\ii\mathcal W_\mu^2}{\sqrt2}.}
 \label{eq:wpm-definition-dp11}
\end{equation}
因为 $(\mathcal W^1)^2+(\mathcal W^2)^2=2\mathcal W^+\mathcal W^-$，第一行成为
\begin{equation*}
 \frac{g^2v^2}{4}\mathcal W_\mu^+\mathcal W^{-\mu}.
\end{equation*}
与几何归一化下的有质量矢量场质量项
\begin{equation*}
 \hbar c\left(\frac{m_Wc}{\hbar}\right)^2
 \mathcal W_\mu^+\mathcal W^{-\mu}
\end{equation*}
比较：
\begin{equation}
 \boxed{m_Wc^2=\frac{gv_E}{2}.}
 \label{eq:sm-w-mass-dp11}
\end{equation}
\end{proposition}

\begin{proposition}[中性质量矩阵的正交对角化]
中性扇区写成矩阵形式：
\begin{equation}
 \Lag_{\rm neutral}^{\rm mass}
 =\frac{v^2}{8}
 \begin{pmatrix}\mathcal W_\mu^3&\mathcal B_\mu\end{pmatrix}
 \begin{pmatrix}g^2&-gg'\\-gg'&g'^2\end{pmatrix}
 \begin{pmatrix}\mathcal W^{3\mu}\\\mathcal B^\mu\end{pmatrix}.
 \label{eq:neutral-mass-matrix-dp11}
\end{equation}
这个矩阵的两个本征值是 $0$ 与 $g^2+g'^2$。零本征值对应未破缺的电磁方向，非零本征值对应有质量中性矢量场。为写出归一化本征矢，定义弱混合角
\begin{equation}
 \boxed{
 \sin\theta_W\equiv s_W=\frac{g'}{\sqrt{g^2+g'^2}},
 \qquad
 \cos\theta_W\equiv c_W=\frac{g}{\sqrt{g^2+g'^2}}.}
 \label{eq:weak-angle-dp11}
\end{equation}
正交变换
\begin{equation}
 \boxed{
 \begin{pmatrix}A_\mu^{\rm geo}\\\mathcal Z_\mu\end{pmatrix}
 =\begin{pmatrix}s_W&c_W\\c_W&-s_W\end{pmatrix}
 \begin{pmatrix}\mathcal W_\mu^3\\\mathcal B_\mu\end{pmatrix}}
 \label{eq:az-mixing-dp11}
\end{equation}
给出质量本征态。
\end{proposition}

\begin{definition}[光子几何归一化与 SI 桥]
标准模型质量基中的光子沿用规范场的几何归一化，记为 $A_\mu^{\rm geo}$；前面 Maxwell/QED 章节的 $A_\mu$ 则是 SI 四势。二者满足
\begin{equation}
 \boxed{
 A_\mu^{\rm geo}
 =\sqrt{\frac{\varepsilon_0c}{\hbar}}\,A_\mu,
 \qquad
 F_{\mu\nu}^{\rm geo}
 =\sqrt{\frac{\varepsilon_0c}{\hbar}}\,F_{\mu\nu},
 \qquad
 g_{\rm em}A_\mu^{\rm geo}=\frac{e}{\hbar}A_\mu.}
 \label{eq:sm-photon-geo-si-bridge-dp75}
\end{equation}
其中 $F_{\mu\nu}^{\rm geo}=\partial_\mu A_\nu^{\rm geo}-\partial_\nu A_\mu^{\rm geo}$，$e>0$ 是元电荷，$g_{\rm em}=e/\sqrt{\varepsilon_0\hbar c}$。因此 $[A_\mu^{\rm geo}]={\rm m^{-1}}$，而 $A_\mu$ 保持 SI 四势量纲；\eqnrefs{eq:sm-photon-geo-si-bridge-dp75} 把标准模型无量纲耦合与 QED 的物理电荷归一化接到同一电磁场上。
\end{definition}

\begin{corollary}[光子、$Z$ 与树级 $\rho$ 参数]
代回 \eqnrefs{eq:neutral-mass-matrix-dp11}：
\begin{equation}
 \boxed{
 m_\gamma=0,
 \qquad
 m_Zc^2=\frac{v_E}{2}\sqrt{g^2+g'^2}
 =\frac{m_Wc^2}{c_W}.}
 \label{eq:sm-z-photon-masses-dp11}
\end{equation}
光子无质量来自 Higgs 真空仍保留一个未破缺规范生成元：对应方向正是中性质量矩阵的零本征矢。由同一质量关系还可定义树级电弱参数
\begin{equation}
 \boxed{
 \rho_{\rm tree}
 \equiv\frac{(m_Wc^2)^2}{(m_Zc^2)^2c_W^2}=1.}
 \label{eq:sm-rho-tree-dp85}
\end{equation}
这里的等号只使用 Higgs 二重态给出的质量关系 $m_Zc^2=m_Wc^2/c_W$；圈修正或不同 Higgs 表示会使精密电弱参数偏离这一树级值。
\end{corollary}

\begin{proposition}[正半定中性质量矩阵的分解]
设 Higgs 真空与正动能归一化按正文固定，质量矩阵在实中性规范场空间上对称。则有以下结论（原式 \eqnrefs{eq:sm-z-photon-masses-dp11}）：
\begin{equation*}
 \boxed{
 m_\gamma=0,
 \qquad
 m_Zc^2=\frac{v_E}{2}\sqrt{g^2+g'^2}
 =\frac{m_Wc^2}{c_W}.}
\end{equation*}
\end{proposition}
\begin{proof}
正文\eqnrefs{eq:sm-gauge-mass-before-diagonal-dp11} 已给出 Higgs 真空诱导的带电与中性规范场二次型。带电扇区直接与有质量矢量场的标准质量项比较；中性扇区则对角化 $2\times2$ 质量矩阵。两步合在一起给出正文\eqnrefs{eq:sm-z-photon-masses-dp11} 中的零质量光子以及 $W$、$Z$ 质量。

对 Higgs 真空，
\begin{equation*}
 D_\mu\langle\Phi\rangle
 =\frac{\ii v}{2\sqrt2}
 \left[g\sigma^I\mathcal W_\mu^I+g'\mathcal B_\mu\right]
 \begin{pmatrix}0\\1\end{pmatrix}.
\end{equation*}
Pauli 矩阵作用在下分量上给出
\begin{equation*}
 \sigma^1\binom01=\binom10,
 \qquad
 \sigma^2\binom01=\binom{-\ii}0,
 \qquad
 \sigma^3\binom01=\binom0{-1}.
\end{equation*}
于是
\begin{equation*}
 D_\mu\langle\Phi\rangle
 =\frac{\ii v}{2\sqrt2}
 \begin{pmatrix}
 g(\mathcal W_\mu^1-\ii\mathcal W_\mu^2)\\
 -g\mathcal W_\mu^3+g'\mathcal B_\mu
 \end{pmatrix},
\end{equation*}
取模平方即得到\eqnrefs{eq:sm-gauge-mass-before-diagonal-dp11}。

带电扇区使用 $\mathcal W^\pm=(\mathcal W^1\mp\ii\mathcal W^2)/\sqrt2$，所以
\begin{equation*}
 \frac{v^2g^2}{8}\left[(\mathcal W^1)^2+(\mathcal W^2)^2\right]
 =\frac{g^2v^2}{4}\mathcal W^+_\mu\mathcal W^{-\mu},
\end{equation*}
与有质量矢量场的质量项比较得到 $m_Wc^2=gv_E/2$。

中性扇区的无量纲矩阵为
\begin{equation*}
 M_0^2=\begin{pmatrix}g^2&-gg'\\-gg'&g'^2\end{pmatrix}.
\end{equation*}
其特征多项式为
\begin{align*}
 \det(M_0^2-\lambda I)
 &=(g^2-\lambda)(g'^2-\lambda)-g^2g'^2\\
 &=\lambda\left[\lambda-(g^2+g'^2)\right].
\end{align*}
因此本征值是 $0$ 与 $g^2+g'^2$；相应的归一化本征矢正是\eqnrefs{eq:az-mixing-dp11} 中的光子方向与 $Z$ 方向。把非零本征值乘回 $v^2/4$ 并换成 $v_E$，得到\eqnrefs{eq:sm-z-photon-masses-dp11}。
\end{proof}


\subsection{电弱相互作用与规范量子化}

\begin{proposition}[电荷生成元来自真空稳定子]
零质量规范场对应一个保持真空不变的生成元，因此要求某个线性组合
\begin{equation*}
 Q=aT^3+bY
\end{equation*}
湮灭 Higgs 真空。Higgs 真空的下分量具有
\begin{equation*}
 T^3=-\frac12,
 \qquad Y=+\frac12.
\end{equation*}
因此
\begin{equation*}
 Q\langle\Phi\rangle
 =\left(-\frac a2+\frac b2\right)\langle\Phi\rangle=0
\end{equation*}
要求 $a=b$。选择通常的归一化 $a=b=1$，得到
\begin{equation}
 \boxed{Q=T^3+Y.}
 \label{eq:gm-nishijima-sm-dp11}
\end{equation}
由 $Q=T^3+Y$ 可逐分量求电荷：
\begin{center}
\begin{tabular}{c|c|c|c}
场分量 & $T_3$ & $Y$ & $Q=T_3+Y$\\
\hline
$u_L$ & $+1/2$ & $+1/6$ & $+2/3$\\
$d_L$ & $-1/2$ & $+1/6$ & $-1/3$\\
$u_R$ & $0$ & $+2/3$ & $+2/3$\\
$d_R$ & $0$ & $-1/3$ & $-1/3$\\
$\nu_L$ & $+1/2$ & $-1/2$ & $0$\\
$e_L$ & $-1/2$ & $-1/2$ & $-1$\\
$e_R$ & $0$ & $-1$ & $-1$\\
$\phi^+$ & $+1/2$ & $+1/2$ & $+1$\\
$\phi^0$ & $-1/2$ & $+1/2$ & $0$
\end{tabular}
\end{center}
因此真空稳定子生成元同时确定观测到的电磁电荷结构；同一个 $Q$ 对左、右手分量给出相同电荷，使对称性破缺后的 Dirac 质量项仍保持 $U(1)_{\rm em}$ 不变。

电荷公式因此直接来自“真空保持不变”的群论条件，并统一生成各多重态分量的电荷。
\end{proposition}

\begin{proposition}[真空稳定子的电荷生成元]
设电弱真空方向与超荷归一化已固定；未破缺生成元必须消去该真空。则有以下结论（原式 \eqnrefs{eq:gm-nishijima-sm-dp11}）：
\begin{equation*}
 \boxed{Q=T^3+Y.}
\end{equation*}
\end{proposition}
\begin{proof}
先不代入标准模型 Higgs 二重态的具体量子数。设一个标量多重态在真空方向
\[
|\Omega\rangle
\]
上取得非零期望值，并且该方向同时是 \(T^3\) 与 \(Y\) 的本征态：
\[
T^3|\Omega\rangle
=
m\,|\Omega\rangle,
\qquad
Y|\Omega\rangle
=
y\,|\Omega\rangle.
\]
这里 \(m\) 与 \(y\) 暂时保留一般值。中性电弱协变导数作用在均匀真空上时只留下
\[
D_\mu\langle\Phi\rangle
=
-\ii
\left(
g\,m\,\mathcal W_\mu^3
+
g' y\,\mathcal B_\mu
\right)
\langle\Phi\rangle.
\]
定义真空模长
\[
\Xi_\Omega\equiv
\langle\Phi\rangle^\dagger\langle\Phi\rangle>0.
\]
由
\(
(D_\mu\Phi)^\dagger D^\mu\Phi
\)
得到的中性规范场二次项严格为
\[
\mathcal L_{\rm mass}^{(0)}
=\Xi_\Omega
\left(
gm\,\mathcal W_\mu^3+g'y\,\mathcal B_\mu
\right)
\left(
gm\,\mathcal W^{3\mu}+g'y\,\mathcal B^\mu
\right).
\]
按
\[
\mathcal L_{\rm mass}^{(0)}
=\frac12
\begin{pmatrix}\mathcal W_\mu^3&\mathcal B_\mu\end{pmatrix}
\mathsf M_0^2
\begin{pmatrix}\mathcal W^{3\mu}\\\mathcal B^\mu\end{pmatrix}
\]
定义质量矩阵，得到精确等式
\[
\mathsf M_0^2
=2\Xi_\Omega
\begin{pmatrix}
g^2m^2 & gg'my\\
gg'my & g'^2y^2
\end{pmatrix}
=2\Xi_\Omega
\begin{pmatrix}gm\\ g'y\end{pmatrix}
\begin{pmatrix}gm & g'y\end{pmatrix}.
\]
因此 $\operatorname{rank}\mathsf M_0^2\le1$。当 $(m,y)\neq(0,0)$ 时，一个归一化零本征向量可取
\[
\mathbf v_0
=\frac1{\sqrt{g'^2y^2+g^2m^2}}
\begin{pmatrix}
g'y\\-gm
\end{pmatrix},
\qquad
\mathsf M_0^2\mathbf v_0=0.
\]
因此，只要真空还保留一个连续 \(U(1)\) 方向，中性规范场质量矩阵就必有一个无质量组合；这个结论先于任何弱混合角的具体定义。

同一个结果也可完全在生成元空间中得到。设未破缺生成元为
\[
Q_{\rm unbroken}
=
aT^3+bY.
\]
真空稳定条件要求
\[
Q_{\rm unbroken}|\Omega\rangle=0,
\]
即
\[
am+by=0.
\]
若 \(y\neq0\)，可把整体归一化选成 \(a=1\)，于是
\[
\boxed{
Q_{\rm unbroken}
=
T^3-\frac{m}{y}Y.
}
\]
因此未破缺生成元由真空在两个内部方向上的量子数决定，而不是事先指定一个电荷公式。

现在才代入标准模型 Higgs 真空。采用正文 convention，
\[
\langle\Phi\rangle
=
\frac1{\sqrt2}
\begin{pmatrix}0\\v_E/c^2\end{pmatrix},
\]
其下分量满足
\[
m=-\frac12,
\qquad
y=+\frac12.
\]
因此
\[
Q_{\rm unbroken}
=
T^3+Y.
\]
这就是正文\eqnrefs{eq:gm-nishijima-sm-dp11} 的电荷生成元。

再看规范场零模。代入 $m=-1/2$、$y=+1/2$ 后，归一化零本征向量为
\[
\mathbf v_0
=\frac1{\sqrt{g^2+g'^2}}
\begin{pmatrix}g'\\g\end{pmatrix}.
\]
定义
\[
s_W
=
\frac{g'}{\sqrt{g^2+g'^2}},
\qquad
c_W
=
\frac{g}{\sqrt{g^2+g'^2}},
\]
归一化零模就是
\[
A_\mu
=
s_W\mathcal W_\mu^3
+
c_W\mathcal B_\mu,
\]
与正文\eqnrefs{eq:az-mixing-dp11} 一致。与它正交的组合
\[
Z_\mu
=
c_W\mathcal W_\mu^3
-
s_W\mathcal B_\mu
\]
沿着真空实际被规范协变导数“看见”的方向，因此获得质量。

最后可以直接核对耦合。中性规范相互作用
\[
gT^3\mathcal W_\mu^3
+
g'Y\mathcal B_\mu
\]
代入
\[
\mathcal W_\mu^3=c_WZ_\mu+s_WA_\mu,
\qquad
\mathcal B_\mu=-s_WZ_\mu+c_WA_\mu
\]
后，光子项为
\[
\left(
gs_WT^3+g'c_WY
\right)A_\mu.
\]
由于
\[
gs_W=g'c_W
\equiv e,
\]
得到
\[
e(T^3+Y)A_\mu
=
eQ_{\rm unbroken}A_\mu.
\]
因此“\(Q\) 湮灭 Higgs 真空”“中性质量矩阵存在零本征矢”“该零模以同一个电荷 \(Q\) 耦合到所有场”是同一个未破缺局域 \(U(1)\) 对称性的三种等价表述。
\end{proof}


\begin{proposition}[光子与 $Z$ 中性流生成元]
对任意场，把中性协变导数中的规范部分
\begin{equation*}
 gT^3\mathcal W_\mu^3+g'Y\mathcal B_\mu
\end{equation*}
用 \eqnrefs{eq:az-mixing-dp11} 反解：
\begin{align*}
 \mathcal W^3_\mu&=s_WA_\mu^{\rm geo}+c_W\mathcal Z_\mu,\\
 \mathcal B_\mu&=c_WA_\mu^{\rm geo}-s_W\mathcal Z_\mu.
\end{align*}
于是
\begin{align*}
 gT^3\mathcal W^3+g'Y\mathcal B
 =&\left(gs_WT^3+g'c_WY\right)A\\
 &+\left(gc_WT^3-g's_WY\right)\mathcal Z.
\end{align*}
由 \eqnrefs{eq:weak-angle-dp11} 定义
\begin{equation}
 \boxed{g_{\rm em}\equiv gs_W=g'c_W.}
 \label{eq:gem-definition-dp11}
\end{equation}
第一项变成
\begin{equation*}
 g_{\rm em}(T^3+Y)A
 =g_{\rm em}QA.
\end{equation*}
第二项先用 $Y=Q-T^3$：
\begin{align*}
 gc_WT^3-g's_WY
 &=gc_WT^3-g's_W(Q-T^3),\\
 &=\left(gc_W+g's_W\right)T^3-g's_WQ,\\
 &=\frac g{c_W}T^3-\frac g{c_W}s_W^2Q.
\end{align*}
所以
\begin{equation}
 \boxed{
 gT^3\mathcal W^3+g'Y\mathcal B
 =g_{\rm em}QA
 +\frac g{c_W}\left(T^3-s_W^2Q\right)\mathcal Z.}
 \label{eq:neutral-current-generator-dp11}
\end{equation}

几何归一化的光子场与 SI 制 Maxwell 场归一化之间的耦合常数满足
\begin{equation}
 \boxed{
 g_{\rm em}=\frac{e}{\sqrt{\varepsilon_0\hbar c}},
 \qquad
 \alpha=\frac{g_{\rm em}^2}{4\pi}
 =\frac{e^2}{4\pi\varepsilon_0\hbar c}.}
 \label{eq:gem-si-charge-dp11}
\end{equation}
从而把无量纲电弱耦合 $g_{\rm em}$ 与 SI 制的正基本电荷 $e>0$ 明确区分。
\end{proposition}

\begin{proposition}[阶梯算符、$W$ 电荷与左手带电流]
带电流来自 $T^1\mathcal W^1+T^2\mathcal W^2$。定义 $T^\pm=T^1\pm\ii T^2$ 后，规范基中的横向弱同位旋耦合可改写为
\begin{equation}
\boxed{
 g(T^1\mathcal W^1+T^2\mathcal W^2)
 =\frac g{\sqrt2}
 \left(T^+\mathcal W^++T^-\mathcal W^-\right).}
\label{eq:sm-charged-generator-dp68}
\end{equation}
$\mathcal W^\pm$ 被称为“带电”弱玻色子并不是由实验名称反推得到的。电磁生成元为 $Q=T^3+Y$，而 $SU(2)$ 阶梯算符满足 $[T^3,T^\pm]=\pm T^\pm$，同时超荷生成元与 $SU(2)$ 生成元对易，所以
\begin{equation}
\boxed{[Q,T^\pm]=\pm T^\pm.}
\label{eq:sm-w-charge-commutator-dp122}
\end{equation}
对易关系本身还可以直接变成有限变换。令 $U_Q(\alpha)=\ee^{\ii\alpha Q}$，则 Baker--Hausdorff 级数逐阶重复同一个对易子，因而
\begin{equation*}
U_Q(\alpha)T^\pm U_Q^{-1}(\alpha)
=\ee^{\pm\ii\alpha}T^\pm.
\end{equation*}
因此规范联络中与 $T^\pm$ 相乘的系数 $\mathcal W^\pm$ 在未破缺电磁群下具有相应的相位权重。对二重态 $(f_u,f_d)^T$，$T^+$ 把下分量变成上分量；若二者电荷满足 $Q_u-Q_d=1$，则费米子流按
\begin{equation*}
\bar f_u\gamma^\mu P_L f_d
\longmapsto
\ee^{-\ii\alpha}\bar f_u\gamma^\mu P_L f_d
\end{equation*}
变换，恰由 $\mathcal W^+\mapsto\ee^{+\ii\alpha}\mathcal W^+$ 抵消。于是 $W^+$ 的正电荷不是名称约定，而是带电流局域项保持 $U(1)_{\rm em}$ 不变所要求的量子数；$W^-$ 同理携带 $-e$。

对二重态 $(f_u,f_d)^T$，带电相互作用拉氏量因此包含
\begin{equation}
 \boxed{
 \Lag_{\rm CC}
 =-\frac{\hbar c\,g}{\sqrt2}
 \left[
 \bar f_u\gamma^\mu P_L f_d\,\mathcal W_\mu^+
 +\bar f_d\gamma^\mu P_L f_u\,\mathcal W_\mu^-
 \right].}
 \label{eq:sm-charged-current-preflavor-dp11}
\end{equation}
只有左手投影算符 $P_L$ 出现，这正是弱相互作用宇称不守恒的代数来源。
\end{proposition}

\begin{proposition}[升降生成元的电荷对易关系]
设弱同位旋生成元满足 $\mathfrak{su}(2)$ 关系，超荷与弱同位旋对易，电荷生成元采用正文归一化。则有以下结论（原式 \eqnrefs{eq:sm-w-charge-commutator-dp122}）：
\begin{equation*}
\boxed{[Q,T^\pm]=\pm T^\pm.}
\end{equation*}
\end{proposition}
\begin{proof}
从 $SU(2)$ 代数
\[
[T^I,T^J]=\ii\epsilon^{IJK}T^K
\]
与 $T^\pm=T^1\pm\ii T^2$ 出发，逐项计算
\begin{align*}
[T^3,T^+]
&=[T^3,T^1]+\ii[T^3,T^2]
=\ii T^2+\ii(-\ii T^1)
=T^+,\\
[T^3,T^-]
&=[T^3,T^1]-\ii[T^3,T^2]
=\ii T^2-\ii(-\ii T^1)
=-T^-.
\end{align*}
直积群 $SU(2)_L\times U(1)_Y$ 的两个生成元作用在不同内部因子上，因此
\[
[Y,T^\pm]=0.
\]
代入 $Q=T^3+Y$ 后，
\begin{align*}
[Q,T^\pm]
&=[T^3+Y,T^\pm]\\
&=[T^3,T^\pm]+[Y,T^\pm]\\
&=\pm T^\pm,
\end{align*}
正好得到正文\eqnrefs{eq:sm-w-charge-commutator-dp122}。
\end{proof}


\begin{definition}[$Z$ 流的手征与矢量轴矢耦合]
中性流对 Dirac 费米子可写成矢量--轴矢形式。由\eqnrefs{eq:neutral-current-generator-dp11}，左右手分量的无量纲 $Z$ 耦合分别为
\begin{equation}
\boxed{
 g_L^f=T_3^f-Q_fs_W^2,
 \qquad
 g_{\rm R}^f=-Q_fs_W^2.}
\label{eq:sm-z-chiral-couplings-dp68}
\end{equation}
定义
\begin{equation*}
 g_V^f=g_L^f+g_{\rm R}^f=T_3^f-2Q_fs_W^2,
 \qquad
 g_A^f=g_L^f-g_{\rm R}^f=T_3^f,
\end{equation*}
于是
\begin{equation}
 \boxed{
 \Lag_Z
 =-\frac{\hbar c\,g}{2c_W}
 \sum_f
 \bar f\gamma^\mu\left(g_V^f-g_A^f\gamma^5\right)f\,\mathcal Z_\mu.}
 \label{eq:z-current-lagrangian-dp11}
\end{equation}
\end{definition}

\begin{figure}[htbp]
\centering
\begin{tikzpicture}[>=Latex,
  basis/.style={draw,rounded corners,align=center,minimum width=31mm,minimum height=15mm},
  rot/.style={draw,rounded corners,align=center,minimum width=34mm,minimum height=15mm},
  line label/.style={fill=white,inner sep=1.2pt,font=\small}]
\node[basis] (gauge) at (-4.0,0) {gauge basis\\$\binom{\mathcal W^3_\mu}{\mathcal B_\mu}$};
\node[rot] (rotation) at (0,0) {Weinberg rotation\\$\begin{pmatrix}s_W&c_W\\ c_W&-s_W\end{pmatrix}$};
\node[basis,minimum width=42mm] (mass) at (4.35,0) {mass basis\\$\binom{A_\mu^{\rm geo}}{\mathcal Z_\mu}$};
\draw[->,thick] (gauge.east)--(rotation.west);
\draw[->,thick] (rotation.east)--(mass.west);
\end{tikzpicture}
\caption{中性电弱规范基到质量本征基的正交旋转。光子方向是 Higgs 真空质量矩阵的零本征矢；$Z$ 方向与其正交，并具有非零质量。}
\label{fig:ew-neutral-mixing-dp11}
\end{figure}

一般 Yang--Mills 场强的非线性项固定三、四规范玻色子自相互作用。电弱规范扇区使用 $SU(2)_L$ 的结构常数 $\epsilon^{IJK}$；先在 $\mathcal W^I$ 实基中写出这些自相互作用，再用电弱混合矩阵旋转到 $W^\pm,A,Z$ 质量本征基。$SU(2)_L$ 场强为
\begin{equation}
 \mathcal W^I_{\mu\nu}
 =\partial_\mu\mathcal W^I_\nu-\partial_\nu\mathcal W^I_\mu
 -g\epsilon^{IJK}\mathcal W^J_\mu\mathcal W^K_\nu,
 \label{eq:sm-su2-fieldstrength-massrule-dp11}
\end{equation}
并定义线性旋度
\begin{equation*}
 w^I_{\mu\nu}\equiv\partial_\mu\mathcal W^I_\nu-\partial_\nu\mathcal W^I_\mu.
\end{equation*}
于是规范动能的二、三、四次项可写成
\begin{align}
 \frac{\Lag_W}{\hbar c}
 ={}&-\frac14w^I_{\mu\nu}w^{I\mu\nu}
 +\frac g2\epsilon^{IJK}w^I_{\mu\nu}\mathcal W^{J\mu}\mathcal W^{K\nu}
 \nonumber\\
 &-\frac{g^2}{4}\epsilon^{IJK}\epsilon^{ILM}
 \mathcal W^J_\mu\mathcal W^K_\nu
 \mathcal W^{L\mu}\mathcal W^{M\nu}.
 \label{eq:sm-su2-kinetic-expanded-dp11}
\end{align}
\eqnrefs{eq:qcd-three-gluon-vertex-dp10} 给出一般 Yang--Mills 三点张量；作 $f^{abc}\to\epsilon^{IJK}$、$g_s\to g$ 的群与耦合替换后，三条物理四动量 $p,q,r$ 全部流入并满足 $p+q+r=0$ 时，$SU(2)$ 三点顶角可直接写成
\begin{equation}
 {
 \Gamma^{abc}_{\mu\nu\rho}(p,q,r)
 =-g\epsilon^{abc}
 \left[
 \eta_{\mu\nu}(p-q)_\rho
 +\eta_{\nu\rho}(q-r)_\mu
 +\eta_{\rho\mu}(r-p)_\nu
 \right].}
 \label{eq:sm-yangmills-three-vertex-dp11}
\end{equation}
其中 $p,q,r$ 均采用\eqnrefs{eq:reduced-si-physical-p-r10} 的物理四动量，量纲为动量；约化图规则已经把 Fourier 变换和外腿归一化中成对出现的 $\hbar$ 因子统一提出。

质量本征基由
\begin{equation*}
 \mathcal W^\pm=\frac{\mathcal W^1\mp\ii\mathcal W^2}{\sqrt2},
 \qquad
 \mathcal W^3=s_WA+c_W\mathcal Z
\end{equation*}
给出。只有 $W^+W^-\mathcal W^3$ 组合具有非零 $SU(2)$ 结构常数，因此中性场旋转立即固定两个三点耦合：
\begin{equation}
 {
 g_{WW\gamma}=g_{\rm em},
 \qquad
 g_{WWZ}=gc_W.}
 \label{eq:sm-wwv-couplings-dp11}
\end{equation}
若 $W^+(p_+,\mu),W^-(p_-,\nu),V(k,\rho)$ 全部流入，$p_++p_-+k=0$，则
\begin{equation}
 {
 \Gamma^{WWV}_{\mu\nu\rho}
 =+\ii g_{WWV}
 \left[
 \eta_{\mu\nu}(p_+-p_-)_\rho
 +\eta_{\nu\rho}(p_--k)_\mu
 +\eta_{\rho\mu}(k-p_+)_\nu
 \right].}
 \label{eq:sm-wwv-vertex-dp11}
\end{equation}
它的洛伦兹张量结构由一般 Yang--Mills 三点顶角继承，复因子与耦合则由带电质量本征基变换固定。

电磁规范不变性要求光子纵向分量不能改变物理振幅。为把这一约束写成可直接核验的代数式，定义自由有质量矢量场的横向逆核部分
\begin{equation*}
 \Pi_{\mu\nu}(p)\equiv
 p^2\eta_{\mu\nu}-p_\mu p_\nu-m_W^2c^2\eta_{\mu\nu}.
\end{equation*}
由\eqnrefs{eq:sm-wwv-vertex-dp11} 可得
\begin{equation}
 {
 k^\rho\Gamma^{WW\gamma}_{\mu\nu\rho}
 =+\ii g_{\rm em}
 \left[\Pi_{\mu\nu}(p_-)-\Pi_{\mu\nu}(p_+)\right].}
 \label{eq:sm-wwgamma-ward-dp41}
\end{equation}
对在壳外腿 $p_\pm^2=m_W^2c^2$ 且 $p_\pm\cdot\varepsilon_\pm=0$，右端夹在物理偏振之间为零，因此把外部光子偏振替换为 $k^\rho$ 后振幅消失。

四规范玻色子顶角同样由\eqnrefs{eq:sm-su2-kinetic-expanded-dp11} 的四次项固定，并通过同一基变换产生 $W^+W^-\gamma\gamma$、$W^+W^-\gamma Z$、$W^+W^-ZZ$ 与 $W^+W^-W^+W^-$。树级不存在 $\gamma\gamma\gamma$、$ZZZ$ 或 $\gamma\gamma Z$：三个纯中性外腿都来自同一个 Cartan 方向 $\mathcal W^3$ 与阿贝尔场 $\mathcal B$，相应结构常数为零。

\phantomsection\label{fig:sm-rule-dependency-dp11}
规范玻色子的自相互作用由同一个 Yang--Mills 场强固定。选定 Higgs 真空并对角化质量矩阵后，规范不变作用量给出质量基中的物理顶角；费米子顶角的味结构由 Yukawa 的左右对角化确定。$R_\xi$ 规范固定同时给出 Goldstone 场与 Faddeev--Popov 鬼场的内部传播子和相应顶角，这些扇区必须与物理内部线一同进入完整振幅。



幺正规范把三个 Goldstone 方向完全吸收到有质量规范场中，因此最适合读取树级质量谱；圈图计算则需要同时控制纵向规范场、Goldstone 场与鬼场的紫外行为。$R_\xi$ 规范通过一个连续参数 $\xi_V$ 把这些非物理自由度放在相同的规范依赖极点上，从而使传播子的高动量幂次计数保持透明。

为保持 SI 量纲一致，把具有 $[\Phi]=\sqrt{\mathrm{J/m}}$ 的标量涨落除以 $\sqrt{\hbar c}$，定义几何量纲场
\begin{equation*}
 \widehat\phi^\pm\equiv\frac{\phi^\pm}{\sqrt{\hbar c}},
 \qquad
 \widehat\chi\equiv\frac{\chi}{\sqrt{\hbar c}},
 \qquad
 \widehat h\equiv\frac{h}{\sqrt{\hbar c}},
\end{equation*}
它们的 SI 单位均为 $\mathrm{m^{-1}}$。同时定义 $W,Z,h$ 的 Compton 逆长度
\begin{equation*}
 \kappa_W\equiv\frac{m_Wc}{\hbar},
 \qquad
 \kappa_Z\equiv\frac{m_Zc}{\hbar},
 \qquad
 \kappa_h\equiv\frac{m_hc}{\hbar}.
\end{equation*}
在笛卡尔参数化中，Higgs 二重态写成
\begin{equation}
 \Phi=
 \begin{pmatrix}
 \phi^+\\[2pt]
 \dfrac{v+h+\ii\chi}{\sqrt2}
 \end{pmatrix},
 \qquad \phi^-=(\phi^+)^*.
 \label{eq:sm-higgs-cartesian-rxi-dp11}
\end{equation}
把它代入 $|D_\mu\Phi|^2-V(\Phi)$ 并只保留真空附近的二次项，得到规范场与标量涨落尚未解耦的二次拉氏量
\begin{align}
 \frac{\Lag_{\Phi}^{(2)}}{\hbar c}
 ={}&\partial_\mu\widehat\phi^-\partial^\mu\widehat\phi^+
 +\frac12(\partial_\mu\widehat h)^2
 +\frac12(\partial_\mu\widehat\chi)^2
 -\frac12\kappa_h^2\widehat h^2
 \nonumber\\
 &+\kappa_W^2\mathcal W_\mu^+\mathcal W^{-\mu}
 +\frac12\kappa_Z^2\mathcal Z_\mu\mathcal Z^\mu
 \nonumber\\
 &+\ii\kappa_W
 \left(
 \mathcal W^-_\mu\partial^\mu\widehat\phi^+
 -\mathcal W^+_\mu\partial^\mu\widehat\phi^-
 \right)
 +\kappa_Z\mathcal Z_\mu\partial^\mu\widehat\chi.
 \label{eq:sm-rxi-quadratic-before-gf-dp11}
\end{align}
最后一行说明直接求逆二次核时，$W_\mu^\pm$ 与 $\phi^\pm$、$Z_\mu$ 与 $\chi$ 会发生双线性混合。$R_\xi$ 规范固定函数取为
\begin{align}
 F^+&=\partial^\mu\mathcal W_\mu^+
 -\ii\xi_W\kappa_W\widehat\phi^+,
 &F^-&=\partial^\mu\mathcal W_\mu^-
 +\ii\xi_W\kappa_W\widehat\phi^-,\\
 F^Z&=\partial^\mu\mathcal Z_\mu
 -\xi_Z\kappa_Z\widehat\chi,
 &F^\gamma&=\partial^\mu A_\mu^{\rm geo},
 \label{eq:sm-rxi-gauge-functions-dp11}
\end{align}
并加入
\begin{equation}
 {
 \Lag_{\rm gf}
 =-\frac{\hbar c}{\xi_W}F^+F^-
 -\frac{\hbar c}{2\xi_Z}(F^Z)^2
 -\frac{\hbar c}{2\xi_\gamma}(F^\gamma)^2.}
 \label{eq:sm-rxi-gf-lagrangian-dp11}
\end{equation}
这些函数中的标量项并非任意附加：其系数恰好由\eqnrefs{eq:sm-rxi-quadratic-before-gf-dp11} 的规范场--Goldstone 混合项固定。展开并分部积分后，所有双线性混合精确相消，同时 Goldstone 场获得规范依赖质量
\begin{equation}
 {
 m_{\phi^\pm}^2=\xi_Wm_W^2,
 \qquad
 m_{\chi}^2=\xi_Zm_Z^2.}
 \label{eq:sm-rxi-goldstone-masses-dp11}
\end{equation}
这里的质量只决定规范固定后的内部传播子极点，不代表新的物理粒子谱线。采用物理四动量 $p^\mu$ 后，相应传播子为
\begin{equation*}
 \Delta_{\phi^\pm}(p)
 =\frac{\ii}{p^2-\xi_Wm_W^2c^2+\ii0^+},
 \qquad
 \Delta_{\chi}(p)
 =\frac{\ii}{p^2-\xi_Zm_Z^2c^2+\ii0^+}.
\end{equation*}

规范固定以后还必须修正路径积分中对同一规范轨道的重复计数。设 $\alpha^b(y)$ 是无穷小规范变换参数，$F^a(x)$ 是所选规范条件。把规范条件沿规范轨道改变多少的线性响应写成
\begin{equation*}
 \mathcal M_{\rm FP}^{ab}(x,y)
 \equiv
 \frac{\delta F^a(x)}{\delta\alpha^b(y)}.
\end{equation*}
这个线性算符称为 Faddeev--Popov 算符。它的行列式正是规范轨道与规范固定截面之间的 Jacobian；在量子场论中通常用一对反交换标量场 $c^a,\bar c^a$ 表示这个行列式，这些辅助场称为 Faddeev--Popov 鬼场。鬼场不是外部可探测粒子，它们只在内部传播子和圈图中补偿非物理规范自由度。

在线性化真空附近，破缺生成元使 Goldstone 场沿规范轨道平移，因此 $\mathcal M_{\rm FP}$ 的二次部分与 Goldstone 二次核具有相同的规范依赖质量尺度。由此得到鬼场质量
\begin{equation}
 {
 m_{c^\pm}^2=\xi_Wm_W^2,
 \qquad
 m_{c^Z}^2=\xi_Zm_Z^2,
 \qquad
 m_{c^\gamma}=0.}
 \label{eq:sm-rxi-ghost-masses-dp11}
\end{equation}
与 QED 不同，电弱 Faddeev--Popov 行列式依赖 Higgs 场和非阿贝尔规范场，所以电弱鬼场在圈图中不脱耦。

对任一有质量规范玻色子 $V=W,Z$，$R_{\xi_V}$ 规范下的约化传播子统一为
\begin{equation}
 {
 D^{(V)}_{\mu\nu}(p)
 =\frac{-\ii}{p^2-m_V^2c^2+\ii0^+}
 \left[
 \eta_{\mu\nu}
 -(1-\xi_V)
 \frac{p_\mu p_\nu}{p^2-\xi_Vm_V^2c^2+\ii0^+}
 \right].}
 \label{eq:sm-rxi-massive-vector-propagator-dp41}
\end{equation}
横向物理极点位于 $p^2=m_V^2c^2$，纵向规范极点位于 $p^2=\xi_Vm_V^2c^2$，后者与相应 Goldstone 和鬼场极点重合。$\xi_V=1$ 给出 Feynman--'t Hooft 规范
\begin{equation*}
 D_{\mu\nu}^{(V)}(p)
 =\frac{-\ii\eta_{\mu\nu}}{p^2-m_V^2c^2+\ii0^+},
\end{equation*}
而 $\xi_V\to\infty$ 把非物理极点推向无穷，并得到幺正规范传播子
\begin{equation*}
 D_{\mu\nu}^{(V)}(p)
 \longrightarrow
 \frac{-\ii}{p^2-m_V^2c^2+\ii0^+}
 \left(\eta_{\mu\nu}-\frac{p_\mu p_\nu}{m_V^2c^2}\right).
\end{equation*}
幺正规范的纵向分子在 $|p^2|\gg m_V^2c^2$ 时不随动量衰减，因此逐图紫外幂次计数较差；有限 $\xi_V$ 的 $R_\xi$ 规范把可重整化结构显式保留在传播子中。物理 $S$ 矩阵最终与 $\xi_V$ 无关，这一结论由 BRST 与 Slavnov--Taylor 恒等式保证，而不是要求每一张 Feynman 图分别与 $\xi_V$ 无关。

\phantomsection\label{fig:sm-rxi-pole-structure-dp11}
$R_\xi$ 规范固定消去规范场与 Goldstone 场的二次混合，并同时确定纵向规范场、Goldstone 场与 Faddeev--Popov 鬼场的传播结构。完整振幅必须对这些规范相关的内部扇区求和；BRST 对称性导出的 Slavnov--Taylor 恒等式约束各项，使物理振幅与规范参数 $\xi$ 无关。











\subsection{Higgs 与规范玻色子的耦合由质量固定}

规范玻色子的质量项和一个 Higgs 外腿的三点耦合来自同一个 $|D_\mu\Phi|^2$ 多项式：质量项取真空因子的二次项，$hVV$ 顶角取其中一次把真空因子替换为径向涨落的项。因此一旦 $m_W,m_Z$ 与真空尺度固定，相应树级 Higgs--规范玻色子耦合也随之固定。

\begin{proposition}[质量项决定的 Higgs--矢量耦合]
在幺正规范中，把 $v\to v+h$ 代入 \eqnrefs{eq:sm-gauge-mass-before-diagonal-dp11}。带电扇区给出
\begin{equation*}
 \frac{g^2(v+h)^2}{4}\mathcal W^+\mathcal W^-
 =\frac{g^2v^2}{4}\mathcal W^+\mathcal W^-
 +\frac{g^2v}{2}h\mathcal W^+\mathcal W^-
 +\frac{g^2}{4}h^2\mathcal W^+\mathcal W^-.
\end{equation*}
使用 $m_Wc^2=gv_E/2$ 与 $v_E=\sqrt{\hbar c}v$，线性 Higgs 耦合可写成
\begin{equation}
 \Lag_{hWW}
 =\frac{2(m_Wc^2)^2}{v_E\sqrt{\hbar c}}
 h\,\mathcal W_\mu^+\mathcal W^{-\mu}.
 \label{eq:hww-coupling-si-dp11}
\end{equation}
中性质量项在 $\mathcal Z_\mu$ 基底中为
\[
 \frac{g_Z^2(v+h)^2}{8}\,
 \mathcal Z_\mu\mathcal Z^\mu,
 \qquad g_Z\equiv\sqrt{g^2+g'^2}.
\]
展开 $(v+h)^2=v^2+2vh+h^2$ 后，线性项的系数为 $g_Z^2v/4$。利用
$ m_Zc^2=g_Zv_E/2$ 与 $v_E=\sqrt{\hbar c}v$，有
\[
 \frac{g_Z^2v}{4}
 =\frac{(m_Zc^2)^2}{v_E\sqrt{\hbar c}},
\]
因此
\begin{equation}
 \Lag_{hZZ}
 =\frac{(m_Zc^2)^2}{v_E\sqrt{\hbar c}}
 h\,\mathcal Z_\mu\mathcal Z^\mu.
 \label{eq:hzz-coupling-si-dp11}
\end{equation}
因此 $hVV$ 耦合与矢量玻色子质量平方成正比，不存在独立于质量生成机制的 Higgs--$W$ 耦合常数。
\end{proposition}

\begin{corollary}[由质量参数表示的 Higgs 自耦合]
Higgs 自耦合同样由 $m_h$ 和 $v_E$ 决定。由 \eqnrefs{eq:sm-higgs-potential-expanded-dp11} 与 $\lambda_H=(m_hc^2)^2/(2v_E^2)$，
\begin{equation}
 V(h)
 =\frac12\frac{m_h^2c^2}{\hbar^2}h^2
 +\frac{(m_hc^2)^2}{2v_E(\hbar c)^{3/2}}h^3
 +\frac{(m_hc^2)^2}{8v_E^2\hbar c}h^4+V_0.
 \label{eq:higgs-self-couplings-mass-form-dp11}
\end{equation}
这给出三 Higgs 与四 Higgs 顶角的全部树级系数。
\end{corollary}

\begin{proposition}[能量形式顶角及相同场的阶乘]
散射与衰变振幅需要使用具有统一外线归一化的场。Higgs 场的几何形式 $\widehat h=h/\sqrt{\hbar c}$ 具有逆长度量纲；
用于下列能量形式约化顶角的标量坐标则是 $h_E=\sqrt{\hbar c}\,h$，
与真空能标 $v_E$ 同量纲。以 $h_E$ 表示 $v_E+h_E$ 并采用相同的外线归一化，
对相互作用作泛函微分，得到规范玻色子质量本征基的能量形式顶角
\begin{align}
 hW^+W^-:&\qquad \ii\frac{2(m_Wc^2)^2}{v_E}\eta_{\mu\nu},\\
 hZZ:&\qquad \ii\frac{2(m_Zc^2)^2}{v_E}\eta_{\mu\nu}.
 \label{eq:sm-higgs-vector-rules-dp43}
\end{align}
两式中的质量平方来自同一个 Higgs 质量生成项，不引入新的独立耦合常数。$hZZ$ 对两个相同 $Z$ 场作泛函微分时产生 $2!$，这解释了其顶角系数与拉氏量\eqnrefs{eq:hzz-coupling-si-dp11} 之间的因子二。费米子质量尚未由 Yukawa 矩阵对角化，因此 $hf\bar f$ 顶角在味结构建立后再与这两式合并。
\end{proposition}

同一个 $|D_\mu\Phi|^2$ 同时给出 $m_W$、$m_Z$、零质量光子本征态以及 $hWW$、$hZZ$ 耦合；未破缺生成元满足 $Q\langle\Phi\rangle=0$，从而给出 $Q=T^3+Y$。因此规范玻色子质量与 Higgs--矢量玻色子耦合并不是彼此独立的参数。

% DP90: hWW/hZZ 顶角只需展开 (v+h)^2 并作外场泛函微分；正文已完整给出质量关系与顶角，不再保留低认知增益的独立 Derivation。

